The formal guidelines for writing a Ph.D.\ thesis mention that the preface
should summarize the general aim of the work. Well, to be short and concise, the
aim of my work was to obtain a Ph.D.\ and get to travel and see new places, and
I believe I can confidently say that I have now reached both goals. However,
this experience was quite challenging and I could never have made it without the
people whom I have met and worked with during these years. Every single one of
them, in different ways, ``nudged'' me to the right direction, allowing me to
eventually succeed. For this reason, I must take enough time and space to
properly express them my gratitude.

Actually, this journey started one year before my Ph.D.\ began, when I did an
Erasmus experience in Leuven during my master's degree. In that occasion, the
very first person I got in touch with was Frank Piessens. To be completely
honest, I had no idea who Frank was at the time, but I saw from the KU Leuven
website that he worked in security and had a friendly profile picture, so I had
decided to contact him asking for some master thesis opportunities. This turned
out to be the very beginning of my adventure. So I must say thank you, Frank,
for being so friendly and for connecting me with all the other people in your
group. Without you, my last four years would not have been possible.

During my Erasmus I met two very important people that made a huge
impact on my subsequent work. The first person that I should mention is of
course JT, my main supervisor during all these years. Thank you for being so
encouraging. Your mere presence during the meetings always made me feel safe,
knowing that you would always make sure that we would not go off-track. Thank
you for that, and I will make sure to bring you some
\emph{nyponsoppa}
every time I travel to Sweden. Next, I wanted to thank Fritz, first my daily
advisor during my master thesis, then my colleague during most of my Ph.D., and
now a very good friend. You have been a source of inspiration for me, and I hope
our paths will cross again in the future.

Of course, I should also mention the people I have met since my Ph.D.\ started.
Moving to Sweden in November 2020 during the pandemic was a real challenge,
especially after seeing the sun setting at 3pm. Thankfully, though, I was in
good company. First and foremost I should thank Christoph, my supervisor at
Ericsson. I have really appreciated our relationship, which was more between
equals rather than supervisor and student. I particularly enjoyed our
brainstorming sessions to design the V2X self-revocation schemes. Second, I
would like to express my sincere gratutide to my managers Hans and Ulf, who
always did everything they could to make me feel as much comfortable as possible
at Ericsson. You really made me feel part of the team. Third, I must mention the
other Ph.D.\ students involved in the 5GhOSTS project: Gerald, Merve and Mykyta.
It was a pleasure meeting and working with you all. In particular, I would like
to thank Merve, with whom I shared the most during these years. Thank you for
being such a good friend. I look forward to your defense and I will make sure to
remind you of your remaining milestones. Finally, thank you to all the other
people with whom I partially shared my working days: The Ericsson Research
Security team in Kista, the DistriNet people and especially the TEE group at KU
Leuven. It was great to share the office spaces and have coffee breaks with such
nice people, and I was honestly impressed by your brilliance and talent.

But it is not only the colleagues that I should mention in this preface. Even
though it was exciting to experience different environments during my Ph.D.,
moving between one place and another took a toll on my personal life. However, I
was lucky enough to find amazing people in Leuven with whom I shared many good
memories. It is not possible to mention everybody here, but I want to especially
thank the people from my football and beach volleyball groups. It has been so
fun to spend time with you all and, especially during this last year, I have
really felt at home. I hope we can keep going for some more time. I must also
mention my friends in Italy, from both Oulx and Torino. I regret that we could
not see each other very often in these last years, but it was always great to
spend time together every time I had the chance to come back. Of course, I need
to particularly mention my family, and especially my parents. I know that it has
not been easy for you to have me far away for such a long time, but thank you
for understanding my needs and supporting my decisions. I really appreciate
that.

Lastly, the formal guidelines also mention that I should thank whomever have
awarded my doctoral scolarship. So, even though I should have asked for your
consent before writing this, I just wanted to say thank you, European Union, for
your support.

\hfill\emph{Gianluca}